\section{Synthetic Experiments}
\label{sec:synth}

We use the released simplified simulator harness to produce $S=12$ synthetic samples across $L=300$ loci at coverages $\{0.1\times,1\times,5\times,30\times\}$, and compare a simplified HAVI-Methyl variant against a FinaleMe-style binary-emission baseline. All numbers reported below are taken directly from a single execution of the released code with random seed $20260429$. The simplified HAVI-Methyl variant uses a Gaussian posterior and empirical-Bayes hierarchical fit on top of the same per-fragment classifier used by the baseline. The full Set Transformer, normalizing-flow posterior, conformal wrapper, mQTL-anchor integration, and full tissue-of-origin head remain planned for the real-data benchmarking phase described in Section~\ref{sec:benchmark}.

\subsection{Per-CpG \texorpdfstring{$\beta$}{beta} recovery}
Table~\ref{tab:recovery} reports per-CpG Pearson correlation and root-mean-square error of predicted versus ground-truth $\beta$ across coverages. The simplified HAVI-Methyl variant improves Pearson correlation at every coverage level, with the largest absolute gains in the low-coverage regime where the hierarchical population prior contributes the most. Figure~\ref{fig:recovery} shows the per-locus scatter of true versus predicted $\beta$ for both methods at each coverage.

\begin{table}[htbp]
\centering
\caption{Per-CpG $\beta$ recovery in the released simplified synthetic harness: Pearson $r$ (higher better) and RMSE (lower better) on $S=12$ samples and $L=300$ loci.}
\label{tab:recovery}
\begin{tabular}{lcccc cccc}
\toprule
& \multicolumn{4}{c}{Pearson $r$} & \multicolumn{4}{c}{RMSE}\\
\cmidrule(lr){2-5}\cmidrule(lr){6-9}
Method & $0.1\times$ & $1\times$ & $5\times$ & $30\times$ & $0.1\times$ & $1\times$ & $5\times$ & $30\times$\\
\midrule
FinaleMe-style HMM & $0.162$ & $0.524$ & $0.768$ & $0.961$ & $0.310$ & $0.312$ & $0.187$ & $0.098$\\
HAVI-Methyl (simplified) & $\mathbf{0.333}$ & $\mathbf{0.782}$ & $\mathbf{0.868}$ & $\mathbf{0.964}$ & $\mathbf{0.296}$ & $\mathbf{0.227}$ & $\mathbf{0.141}$ & $\mathbf{0.096}$\\
\bottomrule
\end{tabular}
\end{table}

\begin{figure}[htbp]
\centering
\includegraphics[width=\textwidth]{figures/recovery_scatter.png}
\caption{Per-locus scatter of true versus predicted $\beta$ at each coverage in the simplified synthetic harness. Both methods approach the diagonal as coverage increases; the simplified HAVI-Methyl variant benefits more at low coverage.}
\label{fig:recovery}
\end{figure}

\subsection{Low-information locus recovery}
The current code uses the bottom 30\% of local methylation variability, computed as the standard deviation of the population $\beta$ track in a $\pm 5$-locus window, as a synthetic proxy for low-information or CpG-poor behavior. On this subset, the FinaleMe-style baseline achieves Pearson correlations of $0.224$, $0.588$, $0.833$, and $0.978$ across the four coverages, while the simplified HAVI-Methyl variant achieves $0.398$, $0.828$, $0.906$, and $0.980$. This proxy should be replaced by an explicit CpG-density or genomic-context stratification in the real-data benchmark.

\subsection{Calibration}
At nominal $90\%$ posterior credible intervals, the simplified HAVI-Methyl variant achieves empirical coverage of $0.471$, $0.597$, $0.671$, and $0.415$ at the four coverages on the full locus set, with mean interval widths of $0.512$, $0.442$, $0.280$, and $0.128$. The empirical coverage undershoots the nominal rate, especially at low and high coverage, which is expected for the simplified Gaussian-posterior variant. The conformal wrapper of Section~\ref{sec:calibration} is the intended calibration mechanism for production evaluation; these raw intervals should not be treated as demonstrated conformal calibration.

\begin{figure}[htbp]
\centering
\includegraphics[width=0.55\textwidth]{figures/calibration.png}
\caption{Calibration curve at $5\times$ coverage under the released simplified-Gaussian variant. The dashed line is ideal coverage. Production evaluation should use the normalizing-flow posterior plus a split-conformal wrapper.}
\label{fig:calibration}
\end{figure}

\subsection{Identifiability stress test}
We synthesize a buffy-coat prior input correlated with disease status by adding a $0.25$ shift on the $\beta$ scale at $25\%$ of loci for half of the $12$ samples. We measure prior attribution as the partial $R^2$ of the prior input in a regression of predicted $\beta$ onto the joint of prior input and ground truth. Without VIB regularization the prior accounts for $5.93\%$ of the predicted signal; adding VIB regularization at $\beta_{\mathrm{VIB}}=0.3$ reduces this to $0.62\%$; adding mQTL anchors on top reduces it further to $0.037\%$. These synthetic reductions are consistent with the intended leakage-control mechanisms, but they are not a proof of real-data identifiability.

\subsection{Tissue-fraction recovery}
Across mixtures of three synthetic tissues, the simplified continuous deconvolution proxy achieves tissue-fraction RMSE of $0.007$ versus $0.129$ for a FinaleMe-derived binarize-and-deconvolve baseline. The result illustrates the cost of binarization in a synthetic setting. It should not be read as a completed evaluation of the full Dirichlet or HDP tissue-of-origin head in Section~\ref{sec:too}.

\subsection{Surrogate objective trajectory}
Figure~\ref{fig:elbo} shows the surrogate objective trajectory for the simplified variant during one run. The plotted quantity is a negative-SSE-style proxy used by the harness, not the exact full-model ELBO. Its decrease is compatible with stable optimization but does not by itself verify Proposition~\ref{prop:svi} for the full flow-based model.

\begin{figure}[htbp]
\centering
\includegraphics[width=0.55\textwidth]{figures/elbo_trajectory.png}
\caption{Surrogate objective trajectory at $30\times$ coverage during synthetic training. The negative-SSE proxy decreases with iteration in the simplified harness; the full-model ELBO should be tracked separately in the planned implementation.}
\label{fig:elbo}
\end{figure}

\subsection{Caveats and ablation pointer}
The synthetic-experiment numbers are single-seed point estimates at reduced scale, not final benchmark statistics. Multi-seed repeats, bootstrap intervals, and full ablations of the Set Transformer encoder, normalizing-flow posterior, IWAE objective, VIB penalty, mQTL anchors, conformal wrapper, and tissue head are planned for Section~\ref{sec:benchmark}. Appendix~\ref{app:extres} reports extended metrics from the same released run and records sensitivity analyses that should be regenerated before publication if they are to be treated as completed measurements.

% ============================================================
